\chapter{The Shape of Seventeen Years}

\bigskip

The first post appeared on January 9, 2008: Moore's law and software bloat, 182 words. A small observation. Nothing viral.

The most recent post, as of this writing, appeared on November 27, 2025: four equivalent formulas for measuring distance in hyperbolic space, with Python code. Still small observations. Still no announcements.

The numbers are simple: 5,233 posts over 6,531 days. The shape of those numbers is the story.

\bookfigure{ch01_activity_heatmap}{Posting activity across 17 years}

\bigskip

\section*{Year by Year}

The shape of the blog changes over time, but the fact of its presence does not. The annual accounting:

\begin{center}
\begin{tabular}{llll}
\toprule
\textbf{Year} & \textbf{Posts} & \textbf{Avg Words} & \textbf{Notes} \\
\midrule
2008 & 410 & 206 & Launch year—highest volume \\
2009 & 366 & 256 &  \\
2010 & 284 & 268 &  \\
2011 & 317 & 222 &  \\
2012 & 357 & 216 &  \\
2013 & 260 & 230 &  \\
2014 & 147 & 188 & Lowest year \\
2015 & 168 & 293 & Recovery begins \\
2016 & 209 & 386 & Posts getting longer \\
2017 & 184 & 383 &  \\
2018 & 266 & 473 & Length nearly doubled from early years \\
2019 & 288 & 446 &  \\
2020 & 288 & 451 & Pandemic year—no change in rhythm \\
2021 & 341 & 451 &  \\
2022 & 362 & 476 & Peak average length \\
2023 & 395 & 435 &  \\
2024 & 248 & 411 &  \\
2025 & 343 & 436 & (partial year through November) \\
\bottomrule
\end{tabular}
\end{center}

The total count obscures the real story.

\bookfigure{ch01_posts_wordcount}{The inverse relationship: fewer posts, more depth}

\bigskip

\section*{The Launch Year}

2008 stands alone. Four hundred and ten posts in a single year—more than one per day on average. Cook posted twice on January 9th, his first day. He posted on January 10th. And 11th. And kept going.

This is the energy of beginning. A new blog, a new outlet, ideas that had been accumulating for years finally finding a place to go. The posts are short—206 words on average, some under 100. They read like notes jotted between other work. A quick observation about software bloat. A pointer to an interesting integral. A link to someone else's insight.

Many new blogs start this way. Few sustain it. The typical pattern is a burst of enthusiasm followed by gradual abandonment. The internet is littered with blogs that have three posts from 2009 and then nothing.

Cook's first year was not a burst. It was the beginning of a practice.

\bigskip

\section*{The Quiet Year}

Then there's 2014. One hundred forty-seven posts—barely one every two and a half days. The average word count dropped to 188, the lowest of any year. Something changed.

The numbers don't say why. The blog offers no announcement, no "I'm taking a break," no explanation of reduced capacity. The posts simply came less frequently for a while, and then more frequently again.

Looking month by month, the reduction is distributed evenly across the year. January: 9 posts. February: 15. March: 11. No single month stands out as empty; no dramatic event visible in the rhythm. Just a year with less to say, or less time to say it, or both.

The surrounding years provide context. 2013: 260 posts. 2014: 147. 2015: 168. The dip begins in 2014, continues into 2015, and then the numbers climb again. By 2016, the blog produces 209 posts—and crucially, those posts are significantly longer. The average jumps from 188 words to 386.

Whatever happened in 2014, the blog emerged from it transformed. The years of many short posts gave way to fewer but deeper investigations. The practice adapted.

\bigskip

\section*{The Trade-Off}

Here is the central pattern of seventeen years: an inverse relationship between quantity and depth.

The early era (2008-2012) produced 1,734 posts averaging 232 words each.
The middle era (2013-2017) produced 968 posts averaging 297 words.
The recent era (2018-2025) produced 2,531 posts averaging 448 words.

The recent era produced over 1.1 million words—about 1.6 times as many as the first two eras combined—despite spanning fewer years. The blog became both more prolific and more substantive. More posts \textit\{and\} longer posts. The assumed trade-off—depth versus frequency—turned out not to apply.

This is what sustained practice looks like. The early years were learning the form. The middle years were consolidation. The recent years are mastery—the ability to produce more, with more substance, without apparent strain.

\bigskip

\section*{The Gaps}

Every blogger takes breaks. The question is how long.

The longest gap in the entire archive is 11 days between posts on July 9 and July 21, 2014. The second-longest gap is 9 days, in May 2014.

No gap reaches two weeks. No sabbatical. No "I'm back after six months." For nearly eighteen years, the blog has never gone silent for more than eleven days.

Most professionals take longer vacations than the blog's longest gap. A two-week holiday would set a new record. A month away would be unprecedented. Whatever else was happening in Cook's life—travel, illness, family, career changes—the blog kept appearing.

The pattern isn't workaholism. The posting frequency is measured, not manic. It's something else: a commitment to showing up, week after week, year after year, regardless of mood or circumstance. The practice survives because it's designed to survive.

\bigskip

\section*{The Pandemic Invariance}

March 2020 upended most patterns. Schools closed, offices emptied, conferences cancelled. The world's rhythms broke.

The blog's rhythms did not.

2020: 288 posts, averaging 451 words. The monthly distribution is unremarkable—22 posts in March (the month everything changed), 21 in April, 23 in May. No spike of anxiety-posting. No drought of distraction. Just another year, indistinguishable from 2019 and 2021.

This flatness is extraordinary. It suggests that the practice serves as a sanctuary—a place where the chaos of the world is ordered into logic.

The blog rarely comments on the pandemic directly. Aside from a handful of posts that use COVID-19 data as statistical examples, it simply continues—another post about probability, another post about Python, another post about number theory. The world outside the blog changed. The blog's function did not.

\bigskip

\section*{What Changed}

Compare the first post to a recent one.

January 9, 2008, "Moore's law and software bloat": A 182-word observation that new software bloats to consume new hardware speed. The post links to one external article. It offers no code, no equations, no deep investigation. It's a thought, shared.

November 27, 2025, "Equal things that don't look equal": A discussion of four different formulas for measuring distance in hyperbolic space, with Python code demonstrating their equivalence, mathematical notation rendered in LaTeX, links to previous posts for context, and a follow-up update about numerical precision. The depth is evident.

Same author. Same platform. Same basic form. But the execution has evolved dramatically. The recent posts assume more context, provide more code, connect more explicitly to the rest of the blog's ecosystem. They read less like notes and more like chapters in an ongoing textbook.

\bigskip

\section*{What Didn't Change}

Still, the core pattern persists. Short, focused posts. One topic per post. A mix of math, computing, statistics, and curiosity. No personal updates, no industry hot takes, no engagement bait. Just: here's something interesting. Here's how it works. Here's some code.

The topics shift over seventeen years—cryptography emerges, PowerShell fades, Python dominates—but the \textit\{kind\} of thinking doesn't change. A working mathematician explaining things to himself, in public, for anyone who happens to be interested.

Cook's own blog provides the frame. In a 2025 post, he observes: "The limiting factor on statistics projects has rarely been my knowledge of statistics. The limiting factors are things like communication, endurance, organization."

Endurance. Organization. The words describe the blog itself.

\bigskip

\section*{Reading the Shape}

Seventeen years of data trace the outline of a practice.

Not a career, though it intersects with one. Not a brand, though others know the name. Not a product, though the writing has value. A practice: something done regularly, for its own sake, getting better over time.

The shape is not linear. It dips in 2014 and rises after. The shape is not stable. Post length doubles over the middle years. The shape is not heroic. Most posts are short, most days uneventful, most numbers unremarkable.

But the shape is complete. Every year filled in. Every month represented. An activity heatmap with no empty regions, no abandoned stretches, no evidence of giving up.

5,233 posts. 1.8 million words. 6,531 days. One voice, present and accounted for.

What does sustained intellectual output look like? This.

\textit{The shape is established. The next question is what fills it—what topics earn that sustained attention, chapter after chapter, year after year.}

